\chapter{Blockchiffren}

\section{Der Übergang}

\begin{merksatz}
\boxtitel{Was Satz 9.3 erzwingt}
Perfekte Sicherheit verlangt einen Schlüssel, der so lang ist wie die Nachricht. Das ist praktisch unbrauchbar.

\begin{center}
\begin{minipage}{0.9\textwidth}
\textbf{Also gibt man sie auf --- und ersetzt sie durch eine schwächere Forderung.}
\end{minipage}
\end{center}

\begin{center}
{\small
\renewcommand{\arraystretch}{1.55}
\begin{tabular}{@{}P{2.6cm}P{6.0cm}@{}}
\toprule
\textbf{Informations\-theoretisch} & Der Geheimtext enthält \kursiv{keine} Information über den Klartext. Gilt gegen jeden Gegner.\\
\textbf{Rechnerisch} & Der Geheimtext enthält Information, aber kein Verfahren mit vertretbarem Aufwand kann sie gewinnen.\\
\bottomrule
\end{tabular}}
\end{center}

\begin{center}
\begin{minipage}{0.88\textwidth}
\kursiv{Der Unterschied ist grundsätzlich. Bei rechnerischer Sicherheit \kursiv{gibt} es die Information --- ein Gegner mit unbegrenzter Rechenzeit findet den Schlüssel stets, einfach durch Durchprobieren. \textbf{Die Sicherheit beruht allein darauf, dass niemand diese Rechenzeit hat.}}
\end{minipage}
\end{center}

\schritt{Die Zahlen}
\begin{center}
{\small
\renewcommand{\arraystretch}{1.48}
\begin{tabular}{@{}P{1.8cm}P{2.2cm}P{3.2cm}@{}}
\toprule
\textbf{Schlüssel} & \textbf{Möglich\-keiten} & \textbf{bei $10^{18}$ je Sekunde}\\
\midrule
$56$ Bit & $7{,}2\cdot10^{16}$ & $0{,}07$ Sekunden\\
$112$ Bit & $5{,}2\cdot10^{33}$ & $1{,}7\cdot10^{8}$ Jahre\\
$128$ Bit & $3{,}4\cdot10^{38}$ & $1{,}1\cdot10^{13}$ Jahre\\
$256$ Bit & $1{,}2\cdot10^{77}$ & $3{,}7\cdot10^{51}$ Jahre\\
\bottomrule
\end{tabular}}
\end{center}

\begin{center}
\begin{minipage}{0.88\textwidth}
\kursiv{Das Weltalter beträgt $1{,}4\cdot10^{10}$ Jahre. Bei $128$ Bit liegt die Suchdauer tausendmal darüber --- und die unterstellte Rate von einer Trillion Schlüsseln je Sekunde ist bereits großzügig. \textbf{Zwischen $56$ und $128$ Bit liegen keine zweieinhalb Größenordnungen, sondern zwanzig.}}
\end{minipage}
\end{center}
\end{merksatz}

\section{Verwirrung und Verwischung}

\begin{merksatz}
\boxtitel{Shannons zwei Entwurfsgrundsätze}
\begin{center}
{\small
\renewcommand{\arraystretch}{1.55}
\begin{tabular}{@{}P{2.4cm}P{6.2cm}@{}}
\toprule
\textbf{Verwirrung} & Der Zusammenhang zwischen Schlüssel und Geheimtext soll so verwickelt sein, dass er sich nicht durch einfache Gleichungen beschreiben lässt.\\
\textbf{Verwischung} & Jedes Bit des Klartextes soll jedes Bit des Geheimtextes beeinflussen --- die Struktur des Klartextes wird über den ganzen Block verteilt.\\
\bottomrule
\end{tabular}}
\end{center}

\begin{center}
\begin{minipage}{0.88\textwidth}
\kursiv{Beide stammen aus Shannons Arbeit von 1949 und sind bis heute die Leitlinien jedes Chiffrenentwurfs. \textbf{Die Verwischung ist die Antwort auf die Häufigkeitsanalyse aus Kapitel~9}: Wenn ein einzelnes Klartextzeichen den ganzen Block verändert, gibt es keine Zeichenhäufigkeiten mehr, die man zählen könnte.}
\end{minipage}
\end{center}
\end{merksatz}

\begin{merksatz}
\boxtitel{Das Lawinenkriterium}
\begin{center}
\begin{minipage}{0.9\textwidth}
\textbf{Ändert man ein Bit der Eingabe, so soll sich im Mittel die Hälfte der Ausgabebits ändern.}
\end{minipage}
\end{center}

\schritt{Warum genau die Hälfte}
Das ist das Verhalten einer zufälligen Abbildung. Jedes Ausgabebit kippt unabhängig mit Wahrscheinlichkeit $\frac12$.

\begin{center}
{\small
\renewcommand{\arraystretch}{1.48}
\begin{tabular}{@{}P{3.4cm}P{1.8cm}P{2.4cm}@{}}
\toprule
& \textbf{Wert} & \\
\midrule
Blocklänge & $128$ & \\
Erwartete Änderung & $64$ & \\
Standard\-abweichung & $5{,}66$ & $\sqrt{128/4}$\\
$\Pr[54\ldots74]$ & $0{,}937$ & \\
\bottomrule
\end{tabular}}
\end{center}

\begin{center}
\begin{minipage}{0.88\textwidth}
\kursiv{Ändern sich deutlich weniger oder deutlich mehr als $64$ Bit, so verhält sich die Chiffre nicht wie eine zufällige Abbildung --- und das ist angreifbar. \textbf{Das Prüfkriterium lautet nicht \enquote{möglichst viel Änderung}, sondern \enquote{genau so viel wie bei Zufall}.} Eine Chiffre, die \kursiv{alle} Bits kippt, wäre ebenso auffällig wie eine, die keines kippt.}
\end{minipage}
\end{center}
\end{merksatz}

\begin{merksatz}
\boxtitel{Der Aufbau in Runden}
Verwirrung und Verwischung erreicht man durch Wiederholung einfacher Bausteine.

\begin{center}
{\small
\renewcommand{\arraystretch}{1.5}
\begin{tabular}{@{}P{2.8cm}P{2.4cm}P{2.6cm}@{}}
\toprule
\textbf{Baustein} & \textbf{Wirkung} & \textbf{Umfang}\\
\midrule
Ersetzung & Verwirrung & wenige Bits\\
Vertauschung & Verwischung & ganzer Block\\
Schlüssel\-addition & bindet den Schlüssel & ganzer Block\\
\bottomrule
\end{tabular}}
\end{center}

\begin{center}
\begin{tikzpicture}[scale=0.95]
\foreach \r in {0,1,2} {
  \begin{scope}[yshift=-\r*1.25cm]
  \foreach \x in {0,1,2,3}
    \draw[tief, line width=0.9pt, fill=tief!10] (\x*1.5,0) rectangle ++(1.0,0.45);
  \node[font=\scriptsize, color=grau, anchor=west] at (6.3,0.22) {Ersetzung};
  \draw[kanted] (0.5,-0.05) -- (2.9,-0.55);
  \draw[kanted] (2.0,-0.05) -- (0.6,-0.55);
  \draw[kanted] (3.5,-0.05) -- (5.0,-0.55);
  \draw[kanted] (5.0,-0.05) -- (3.6,-0.55);
  \end{scope}
}
\node[font=\scriptsize, color=grau, anchor=west] at (6.3,-0.45) {Vertauschung};
\node[font=\scriptsize\itshape, color=grau] at (3.0,-4.2) {mehrere Runden};
\end{tikzpicture}
\end{center}

\begin{center}
\begin{minipage}{0.88\textwidth}
\kursiv{Nach wenigen Runden ist jedes Ausgabebit von jedem Eingabebit abhängig --- das nennt man \Begriff{vollständige Verwischung}. Die tatsächliche Rundenzahl liegt deutlich darüber: Der Advanced Encryption Standard erreicht vollständige Verwischung nach zwei Runden und benutzt zehn bis vierzehn. \textbf{Die überzähligen Runden sind der Sicherheitsabstand}, und der ist es, der ein Verfahren jahrzehntelang trägt.}
\end{minipage}
\end{center}
\end{merksatz}

\section{Betriebsarten}

\begin{merksatz}
\boxtitel{Eine Blockchiffre allein genügt nicht}
Sie verschlüsselt einen Block fester Länge --- etwa $128$ Bit. Nachrichten sind länger.

\begin{center}
{\small
\renewcommand{\arraystretch}{1.55}
\begin{tabular}{@{}P{1.8cm}P{3.0cm}P{3.6cm}@{}}
\toprule
\textbf{Art} & \textbf{Vorschrift} & \textbf{Eigenschaft}\\
\midrule
ECB & $C_i = E_K(M_i)$ & \textcolor{akzent}{unbrauchbar}\\
CBC & $C_i = E_K(M_i\oplus C_{i-1})$ & braucht Startwert\\
CTR & $C_i = M_i\oplus E_K(N\Vert i)$ & Stromchiffre\\
\bottomrule
\end{tabular}}
\end{center}
\end{merksatz}

\begin{warnung}
\boxtitelr{Warum ECB unbrauchbar ist}
\begin{center}
\begin{minipage}{0.9\textwidth}
\textbf{Gleiche Klartextblöcke ergeben gleiche Geheimtextblöcke.}
\end{minipage}
\end{center}

\schritt{Die Folge}
Der Geheimtext verrät, welche Blöcke übereinstimmen --- also die gesamte Wiederholungsstruktur der Nachricht.

\begin{center}
\begin{minipage}{0.88\textwidth}
\kursiv{Bei einem Bild bleiben die Umrisse sichtbar: Flächen gleicher Farbe werden zu Flächen gleicher Geheimtextblöcke. Das bekannteste Anschauungsbeispiel --- ein verschlüsseltes Pinguinbild, dessen Umriss erkennbar bleibt --- ist seit Jahren in jedem Lehrbuch.}
\end{minipage}
\end{center}

\schritt{In der Sprache von Kapitel 4}
\begin{center}
\begin{minipage}{0.9\textwidth}
\textbf{$I(M;C) > 0$, und zwar erheblich.} ECB verletzt die Grundforderung, ohne dass man irgendetwas brechen müsste.
\end{minipage}
\end{center}

\schritt{Bei Befunddaten}
\begin{center}
\begin{minipage}{0.88\textwidth}
\kursiv{Eine Datenbank mit einem Eintrag je Patient, blockweise verschlüsselt: Alle Patienten mit derselben Diagnose haben denselben Geheimtextblock. Man kennt die Diagnose nicht --- aber man kann die Patienten gruppieren und die Gruppengrößen mit der bekannten Häufigkeitsverteilung aus Kapitel~1 abgleichen. \textbf{Damit ist die häufigste Gruppe zugeordnet, und die zweithäufigste auch.} Es ist die Häufigkeitsanalyse aus Kapitel~9, auf Blöcke übertragen.}
\end{minipage}
\end{center}
\end{warnung}

\begin{merksatz}
\boxtitel{Die Abhilfe: Zufall hinzufügen}
\begin{center}
{\small
\renewcommand{\arraystretch}{1.5}
\begin{tabular}{@{}P{2.4cm}P{6.2cm}@{}}
\toprule
\textbf{CBC} & Ein zufälliger Startwert wird dem ersten Block zugemischt; jeder weitere Block erhält den vorangegangenen Geheimtext.\\
\textbf{CTR} & Ein Zähler wird verschlüsselt und liefert einen Schlüsselstrom; die Blockchiffre wird nie auf den Klartext angewandt.\\
\bottomrule
\end{tabular}}
\end{center}

\begin{center}
\begin{minipage}{0.9\textwidth}
\textbf{In beiden Fällen liefern gleiche Klartextblöcke verschiedene Geheimtextblöcke.} Dieselbe Nachricht zweimal verschlüsselt ergibt zwei völlig verschiedene Geheimtexte.
\end{minipage}
\end{center}

\schritt{Der Vergleich}
\begin{center}
{\small
\renewcommand{\arraystretch}{1.48}
\begin{tabular}{@{}P{2.2cm}P{2.0cm}P{2.0cm}P{2.0cm}@{}}
\toprule
& \textbf{ECB} & \textbf{CBC} & \textbf{CTR}\\
\midrule
Parallel & \textcolor{gruen}{ja} & \textcolor{akzent}{nein} & \textcolor{gruen}{ja}\\
Wahlfreier Zugriff & \textcolor{gruen}{ja} & teilweise & \textcolor{gruen}{ja}\\
Auffüllen nötig & ja & ja & \textcolor{gruen}{nein}\\
Sicher & \textcolor{akzent}{nein} & \textcolor{gruen}{ja} & \textcolor{gruen}{ja}\\
\bottomrule
\end{tabular}}
\end{center}

\begin{center}
\begin{minipage}{0.88\textwidth}
\kursiv{CTR verwandelt die Blockchiffre in eine Stromchiffre --- und damit in etwas, das dem Einmalschlüssel aus Kapitel~9 formal gleicht. \textbf{Der Unterschied ist, dass der Schlüsselstrom nicht zufällig, sondern nur \kursiv{nicht unterscheidbar} von zufällig ist.} Damit erbt CTR alle Eigenschaften des Einmalschlüssels --- einschließlich der unangenehmen.}
\end{minipage}
\end{center}
\end{merksatz}

\section{Der Bruchfall}

\begin{warnung}
\boxtitelr{Der Zählerwert darf sich nie wiederholen}
Wird bei CTR derselbe Zählerwert zweimal mit demselben Schlüssel verwendet, so entsteht derselbe Schlüsselstrom:
\[ C_1 \oplus C_2 = M_1 \oplus M_2 \]

\begin{center}
\begin{minipage}{0.9\textwidth}
\textbf{Das ist exakt der Venona-Fehler aus Kapitel~9} --- nur dass er hier nicht durch doppelt gedrucktes Papier entsteht, sondern durch einen Programmfehler.
\end{minipage}
\end{center}

\begin{center}
\begin{minipage}{0.88\textwidth}
\kursiv{Häufige Ursachen: ein Zähler, der nach einem Neustart wieder bei null beginnt; ein zufällig gewählter Startwert, der zu kurz ist; ein auf mehrere Geräte kopierter Schlüssel ohne abgestimmte Zählerbereiche. \textbf{Die Mathematik ist einwandfrei; die Umsetzung ist die Schwachstelle} --- und das ist die durchgehende Erfahrung dieses Fachs.}
\end{minipage}
\end{center}
\end{warnung}

\begin{warnung}
\boxtitelr{Die Geburtstagsschranke}
Bei einer Blocklänge von $n$ Bit treten nach etwa $2^{n/2}$ Blöcken Wiederholungen auf --- rein aus Zufall.

\begin{center}
{\small
\renewcommand{\arraystretch}{1.48}
\begin{tabular}{@{}P{2.2cm}P{1.8cm}P{2.2cm}P{2.2cm}@{}}
\toprule
\textbf{Chiffre} & $n$ & \textbf{Blöcke} & \textbf{Datenmenge}\\
\midrule
3DES & $64$ & $2^{32}$ & $\mathbf{34}$ \textbf{GB}\\
AES & $128$ & $2^{64}$ & $295$ Exabyte\\
\bottomrule
\end{tabular}}
\end{center}

\begin{center}
\begin{minipage}{0.9\textwidth}
\textbf{Bei $64$ Bit Blocklänge genügen $34$ Gigabyte, um eine Kollision wahrscheinlich zu machen.}
\end{minipage}
\end{center}

\begin{center}
\begin{minipage}{0.88\textwidth}
\kursiv{Bei CBC verrät eine Kollision $C_i = C_j$ unmittelbar $M_i \oplus M_j$. Genau darauf beruht ein 2016 veröffentlichter Angriff auf Verbindungen, die noch 3DES verwendeten. \textbf{Das Verfahren war nicht gebrochen --- die Blocklänge war zu kurz.} Bei einer Dauerverbindung sind $34$ Gigabyte binnen Stunden erreicht.}
\end{minipage}
\end{center}

\schritt{Die Lehre}
\begin{center}
\begin{minipage}{0.9\textwidth}
\textbf{Schlüssellänge und Blocklänge sind verschiedene Größen.} $3$DES hat $112$ Bit Schlüssel und ist dennoch begrenzt --- durch die $64$ Bit Blocklänge.
\end{minipage}
\end{center}
\end{warnung}

\begin{warnung}
\boxtitelr{Und drei weitere Fallen}
\begin{center}
{\small
\renewcommand{\arraystretch}{1.5}
\begin{tabular}{@{}P{2.4cm}P{6.2cm}@{}}
\toprule
\textbf{Keine Echtheit} & Auch CBC und CTR sind formbar. Verschlüsselung allein schützt nicht vor gezielter Veränderung.\\
\textbf{Seitenkanäle} & Laufzeit, Stromverbrauch und Zwischenspeicherverhalten können den Schlüssel verraten, obwohl die Mathematik stimmt.\\
\textbf{Auffüllen} & Die Art, wie unvollständige Blöcke ergänzt werden, hat wiederholt zu Angriffen geführt.\\
\bottomrule
\end{tabular}}
\end{center}

\begin{center}
\begin{minipage}{0.88\textwidth}
\kursiv{Die zweite Zeile ist die unangenehmste, weil sie außerhalb des mathematischen Modells liegt. Ein Verfahren, dessen Rechenzeit vom Schlüssel abhängt, verrät ihn --- und die frühen Umsetzungen des Advanced Encryption Standard taten genau das, weil die Ersetzungstabelle im Zwischenspeicher schlüsselabhängig durchlaufen wurde. \textbf{Das Modell dieses Kapitels kennt nur Eingabe und Ausgabe; die Wirklichkeit hat mehr Kanäle.}}
\end{minipage}
\end{center}
\end{warnung}

\section{Übungsaufgaben}

\begin{aufgabe}
\boxtitel{Aufgabe 10.1 --- Der Lawineneffekt}
Eine Chiffre mit $64$ Bit Blocklänge wird geprüft: Bei Änderung eines Eingabebits ändern sich $18$ Ausgabebits. Ist das auffällig?

\loesung
\schritt{Schritt 1 --- Die Erwartung}
Bei einer zufälligen Abbildung kippt jedes Bit unabhängig mit $\frac12$:
\[ \mathbb{E} = 64\cdot0{,}5 = 32, \qquad \sigma = \sqrt{64\cdot0{,}25} = 4 \]

\schritt{Schritt 2 --- Der Abstand}
\[ \frac{18-32}{4} = -3{,}5 \ \text{Standardabweichungen} \]

\schritt{Schritt 3 --- Die Wahrscheinlichkeit}
Nach der Normalverteilungsnäherung (Band 12):
\[ \Pr[X \le 18] \approx \Pr[Z \le -3{,}5] = 2{,}3\cdot10^{-4} \]

\begin{center}
\begin{minipage}{0.9\textwidth}
\textbf{Ein solches Ergebnis tritt bei einer guten Chiffre etwa einmal unter viertausend Versuchen auf --- das ist auffällig.}
\end{minipage}
\end{center}

\schritt{Schritt 4 --- Die Einschränkung}
Prüft man alle $64$ Eingabebits, so hat man $64$ Versuche. Die Wahrscheinlichkeit, dass \kursiv{mindestens einer} so weit abweicht:
\[ 1-(1-2{,}3\cdot10^{-4})^{64} = 1{,}5\,\% \]

\begin{center}
\begin{minipage}{0.88\textwidth}
\kursiv{Das ist das Problem des mehrfachen Testens aus Band 17. \textbf{Ein einzelner auffälliger Wert unter vielen Prüfungen ist noch kein Befund} --- man müsste ihn an weiteren Eingaben bestätigen.}
\end{minipage}
\end{center}

\schritt{Schritt 5 --- Wann es ein Befund wäre}
Träte die Abweichung bei \kursiv{allen} Eingabebits auf, so wäre das systematisch:
\[ (2{,}3\cdot10^{-4})^{64} \approx 10^{-234} \]

\begin{center}
\begin{minipage}{0.88\textwidth}
\kursiv{Dann läge kein Zufall vor, sondern eine Struktur --- und Strukturen sind der Ansatzpunkt jedes Angriffs. \textbf{Die Prüfung des Lawinenkriteriums ist ein statistischer Test wie jeder andere}, mit denselben Fallstricken.}
\end{minipage}
\end{center}
\end{aufgabe}

\begin{aufgabe}
\boxtitel{Aufgabe 10.2 --- Betriebsarten}
Eine Befunddatenbank enthält $50\,000$ Einträge mit je einem Diagnosecode. Die Häufigkeiten sind die aus Kapitel~1. Welche Auskunft gibt ein ECB-verschlüsselter Auszug?

\loesung
\schritt{Schritt 1 --- Was sichtbar wird}
Jeder Diagnosecode ergibt stets denselben Geheimtextblock. Der Angreifer sieht acht verschiedene Blöcke mit ihren Häufigkeiten.

\schritt{Schritt 2 --- Die erwarteten Anzahlen}
\begin{center}
{\small
\renewcommand{\arraystretch}{1.48}
\begin{tabular}{@{}P{1.8cm}P{1.8cm}P{2.2cm}P{2.0cm}@{}}
\toprule
\textbf{Gruppe} & \textbf{Anteil} & \textbf{erwartet} & $\sigma$\\
\midrule
$A$, $B$ & $0{,}2500$ & $12\,500$ & $97$\\
$C$, $D$ & $0{,}1250$ & $6250$ & $74$\\
$E$--$H$ & $0{,}0625$ & $3125$ & $54$\\
\bottomrule
\end{tabular}}
\end{center}

\kursiv{Zur Standardabweichung:} $\sqrt{50\,000\cdot0{,}25\cdot0{,}75} = 96{,}8$.

\schritt{Schritt 3 --- Die Zuordnung}
Die drei Häufigkeitsstufen sind durch je etwa sechzig Standardabweichungen getrennt.

\begin{center}
\begin{minipage}{0.9\textwidth}
\textbf{Die Zuordnung der Stufen ist praktisch sicher.} Innerhalb einer Stufe --- etwa zwischen $A$ und $B$ --- gibt die Häufigkeit nichts her.
\end{minipage}
\end{center}

\schritt{Schritt 4 --- Was der Angreifer erfährt}
\begin{center}
{\small
\renewcommand{\arraystretch}{1.48}
\begin{tabular}{@{}P{2.6cm}P{2.4cm}P{2.6cm}@{}}
\toprule
\textbf{Auskunft} & \textbf{erhalten?} & \\
\midrule
Stufe je Eintrag & \textcolor{akzent}{ja} & drei Klassen\\
Genaue Gruppe & teilweise & bei $C$/$D$ zweideutig\\
Gleichheit zweier Einträge & \textcolor{akzent}{ja} & vollständig\\
\bottomrule
\end{tabular}}
\end{center}

\schritt{Schritt 5 --- Der Informationsgehalt}
Der Angreifer erfährt die Stufe. Die verbleibende Ungewissheit:
\begin{center}
\renewcommand{\arraystretch}{1.46}
\begin{tabular}{@{}P{2.2cm}P{2.2cm}P{2.6cm}@{}}
\toprule
\textbf{Stufe} & \textbf{Anteil} & \textbf{Restentropie}\\
\midrule
$\{A,B\}$ & $0{,}50$ & $1$ Bit\\
$\{C,D\}$ & $0{,}25$ & $1$ Bit\\
$\{E,\ldots,H\}$ & $0{,}25$ & $2$ Bit\\
\bottomrule
\end{tabular}
\end{center}

\[ \Ent(M\mid C) = 0{,}5\cdot1 + 0{,}25\cdot1 + 0{,}25\cdot2 = 1{,}25 \ \text{Bit} \]
\[ I(M;C) = \Ent(M) - \Ent(M\mid C) = 2{,}75 - 1{,}25 = \mathbf{1{,}50} \ \text{Bit} \]

\begin{center}
\begin{minipage}{0.9\textwidth}
\textbf{Mehr als die Hälfte der Information ist durchgesickert --- ohne dass die Chiffre gebrochen wurde.}
\end{minipage}
\end{center}

\schritt{Schritt 6 --- Die Abhilfe}
Bei CBC oder CTR mit wechselndem Startwert ist $I(M;C) = 0$ im Rahmen der rechnerischen Ununterscheidbarkeit.

\begin{center}
\begin{minipage}{0.88\textwidth}
\kursiv{Für eine \kursiv{durchsuchbare} Datenbank ist das allerdings unbequem: Man kann dann nicht mehr nach dem Geheimtext suchen. Verfahren, die beides leisten --- Suchbarkeit und Geheimhaltung ---, geben notwendig etwas preis; die Frage ist nur, wie viel. \textbf{Das Spannungsfeld zwischen Nutzbarkeit und Geheimhaltung lässt sich verschieben, nicht auflösen.}}
\end{minipage}
\end{center}
\end{aufgabe}

\begin{aufgabe}
\boxtitel{Aufgabe 10.3 --- Zwei Schlüssel sind nicht doppelt so gut}
Zeigen Sie, dass eine zweifache Anwendung einer Chiffre mit $k$-Bit-Schlüssel nicht die Sicherheit von $2k$ Bit erreicht.

\loesung
\schritt{Schritt 1 --- Das Verfahren}
\[ C = E_{K_2}\bigl(E_{K_1}(M)\bigr) \]
Der Gesamtschlüssel hat $2k$ Bit, der Raum umfasst $2^{2k}$ Möglichkeiten.

\schritt{Schritt 2 --- Der Angriff}
Der Angreifer kennt ein Paar $(M,C)$. Er nutzt
\[ E_{K_1}(M) = D_{K_2}(C) . \]

\begin{center}
{\small
\renewcommand{\arraystretch}{1.5}
\begin{tabular}{@{}P{1.0cm}P{7.4cm}@{}}
\toprule
$1$ & Berechne $E_{K}(M)$ für alle $2^{k}$ Schlüssel und speichere die Ergebnisse.\\
$2$ & Berechne $D_{K}(C)$ für alle $2^{k}$ Schlüssel.\\
$3$ & Suche Übereinstimmungen. Jede liefert einen Kandidaten $(K_1,K_2)$.\\
\bottomrule
\end{tabular}}
\end{center}

\schritt{Schritt 3 --- Der Aufwand}
\[ 2^{k} + 2^{k} = 2^{\,k+1} \ \text{Verschlüsselungen} \]
statt $2^{2k}$.

\begin{center}
{\small
\renewcommand{\arraystretch}{1.48}
\begin{tabular}{@{}P{1.6cm}P{2.0cm}P{2.2cm}P{2.0cm}@{}}
\toprule
$k$ & \textbf{erhofft} & \textbf{tatsächlich} & \textbf{Speicher}\\
\midrule
$56$ & $2^{112}$ & $2^{57}$ & $2^{56}$ Einträge\\
$128$ & $2^{256}$ & $2^{129}$ & $2^{128}$\\
\bottomrule
\end{tabular}}
\end{center}

\begin{center}
\begin{minipage}{0.9\textwidth}
\textbf{Die zweifache Anwendung bringt genau ein Bit zusätzliche Sicherheit --- nicht $k$.}
\end{minipage}
\end{center}

\schritt{Schritt 4 --- Der Vorbehalt}
Der Angriff braucht $2^{k}$ Speicherplätze. Bei $k=56$ sind das $7{,}2\cdot10^{16}$ Einträge zu je $8$ Byte --- rund $576$ Petabyte.

\begin{center}
\begin{minipage}{0.88\textwidth}
\kursiv{Das ist nicht praktikabel, aber es ist der Punkt: \textbf{Ein Verfahren, dessen Sicherheit allein am Speicherbedarf eines Angriffs hängt, ist nicht solide entworfen.} Es gibt zudem Varianten, die Speicher gegen Rechenzeit tauschen und dabei deutlich unter $2^{2k}$ bleiben.}
\end{minipage}
\end{center}

\schritt{Schritt 5 --- Die Folgerung}
Deshalb wurde dreifache statt zweifacher Anwendung genormt:
\[ C = E_{K_3}\bigl(D_{K_2}(E_{K_1}(M))\bigr) \]

\begin{center}
\renewcommand{\arraystretch}{1.48}
\begin{tabular}{@{}P{2.4cm}P{2.2cm}P{2.6cm}@{}}
\toprule
& \textbf{Schlüssel} & \textbf{Sicherheit}\\
\midrule
einfach & $56$ & $2^{56}$\\
zweifach & $112$ & $2^{57}$\\
dreifach & $168$ & $\approx 2^{112}$\\
\bottomrule
\end{tabular}
\end{center}

\begin{center}
\begin{minipage}{0.88\textwidth}
\kursiv{Auch die dreifache Anwendung erreicht nicht $2^{168}$, sondern nur etwa $2^{112}$ --- derselbe Angriff wirkt eine Stufe höher. \textbf{Die mittlere Entschlüsselung statt Verschlüsselung ist kein Sicherheitsmerkmal}, sondern erlaubt es, mit $K_1=K_2=K_3$ die einfache Chiffre nachzubilden und so ältere Geräte weiterzuverwenden.}
\end{minipage}
\end{center}
\end{aufgabe}

\begin{geschichte}
\boxtitelg{Zwei Normen und ein Wettbewerb}
\Person{Horst Feistel} (1915--1990), aus Deutschland in die Vereinigten Staaten emigriert, entwickelte bei IBM in den 1970er Jahren die Struktur, die seinen Namen trägt.

\begin{center}
\begin{minipage}{0.9\textwidth}
\kursiv{Ihr Kunstgriff: Man teilt den Block in zwei Hälften und verändert in jeder Runde nur eine davon. Dadurch ist die Umkehrung stets möglich --- \kursiv{unabhängig davon}, wie kompliziert die verwendete Funktion ist. Man kann also eine beliebig verwickelte Funktion einbauen, ohne sich um ihre Umkehrbarkeit zu sorgen.}
\end{minipage}
\end{center}

\schritt{Der Data Encryption Standard}
1977 als Norm verabschiedet, mit $56$ Bit Schlüssel.

\begin{center}
\begin{minipage}{0.9\textwidth}
\kursiv{IBM hatte ursprünglich $64$ Bit vorgeschlagen; die Kürzung auf $56$ geschah auf Betreiben der NSA und war von Anfang an umstritten. 1998 baute die Electronic Frontier Foundation für rund $250\,000$ Dollar eine Maschine, die einen Schlüssel in gut zwei Tagen fand --- und beendete damit die Debatte.}
\end{minipage}
\end{center}

\schritt{Die andere Seite}
Zugleich hatte die NSA die Ersetzungstabellen verändert. Das erregte Verdacht.

\begin{center}
\begin{minipage}{0.9\textwidth}
\kursiv{1990 fanden \Person{Eli Biham} und \Person{Adi Shamir} die \Begriff{differentielle Kryptanalyse}. Es zeigte sich, dass die veränderten Tabellen gerade dagegen gehärtet waren --- IBM und die NSA hatten das Verfahren also fünfzehn Jahre vor seiner öffentlichen Entdeckung gekannt und stillschweigend abgewehrt. \textbf{Der Verdacht war berechtigt und die Änderung war eine Verbesserung} --- beides zugleich.}
\end{minipage}
\end{center}

\schritt{Der Wettbewerb}
1997 schrieb das amerikanische Normungsinstitut einen offenen Wettbewerb aus. Fünfzehn Vorschläge wurden eingereicht, öffentlich analysiert und über drei Jahre auf zwei Konferenzen diskutiert.

\begin{center}
\begin{minipage}{0.9\textwidth}
\kursiv{2000 gewann Rijndael der Belgier \Person{Joan Daemen} und \Person{Vincent Rijmen}. Das Verfahren verzichtet auf die Feistel-Struktur und beruht auf Rechnung in einem endlichen Körper --- Band 6, Teil~V. \textbf{Der offene Wettbewerb gilt heute als Vorbild}: Er hat mehr Vertrauen geschaffen als jede Behördenentscheidung es gekonnt hätte, und dasselbe Verfahren wurde später für Hashfunktionen wiederholt.}
\end{minipage}
\end{center}
\end{geschichte}

\begin{forschung}
\boxtitelp{Wo Blockchiffren heute stehen}
\begin{center}
{\small
\renewcommand{\arraystretch}{1.45}
\begin{tabular}{@{}P{2.8cm}P{5.4cm}@{}}
\toprule
\textbf{Stand} & Der Advanced Encryption Standard ist seit über zwanzig Jahren ungebrochen; die besten bekannten Angriffe sparen wenige Bit.\\
\textbf{Verbundene Betriebs\-arten} & Verschlüsselung und Beglaubigung in einem Durchgang --- heute der Normalfall.\\
\textbf{Leichtgewichtig} & Chiffren für Geräte mit wenigen Kilobyte Speicher.\\
\bottomrule
\end{tabular}}
\end{center}

\schritt{Zur ersten Zeile}
Der beste bekannte Angriff auf die Vollrundenfassung mit $128$ Bit senkt den Aufwand von $2^{128}$ auf etwa $2^{126}$.

\begin{center}
\begin{minipage}{0.88\textwidth}
\kursiv{Das ist praktisch bedeutungslos --- $2^{126}$ ist so unerreichbar wie $2^{128}$. Bemerkenswert ist etwas anderes: \textbf{In über zwanzig Jahren intensiver öffentlicher Analyse wurde nichts Wesentliches gefunden.} Für ein Verfahren ohne Sicherheitsbeweis ist das die einzige verfügbare Evidenz --- und sie ist schwächer, als man sich wünschte.}
\end{minipage}
\end{center}

\schritt{Zur zweiten Zeile}
Nach dem Bruchfall dieses Kapitels ist Verschlüsselung ohne Beglaubigung gefährlich.

\begin{center}
\begin{minipage}{0.88\textwidth}
\kursiv{Die Erfahrung hat gezeigt, dass Entwickler beides nur selten richtig kombinieren --- die Reihenfolge ist entscheidend, und die falsche hat wiederholt zu Angriffen geführt. \textbf{Die Lösung bestand darin, die Wahl abzuschaffen}: Moderne Betriebsarten leisten beides und lassen sich gar nicht erst falsch zusammensetzen. Kapitel~11 behandelt die Beglaubigung.}
\end{minipage}
\end{center}
\end{forschung}

\begin{merksatz}
\boxtitel{Was Sie jetzt können}
\begin{center}
\renewcommand{\arraystretch}{1.5}
\begin{tabular}{@{}P{3.0cm}P{5.6cm}@{}}
\toprule
\textbf{Rechnerisch} & Information ist da, aber unerreichbar\\
\textbf{Zwei Grundsätze} & Verwirrung und Verwischung\\
\textbf{Lawine} & genau die Hälfte, nicht mehr\\
\textbf{Betriebsarten} & ECB verrät die Struktur\\
\textbf{Zwei Schranken} & Schlüssellänge und Blocklänge\\
\bottomrule
\end{tabular}
\end{center}

\begin{center}
\begin{minipage}{0.88\textwidth}
\kursiv{Im nächsten Kapitel geht es um die Eigenschaft, die alle bisherigen Verfahren nicht hatten: nachzuweisen, dass eine Nachricht unverändert und vom richtigen Absender ist.}
\end{minipage}
\end{center}
\end{merksatz}
